\textbf{Continued action updating can retain tissue without sustaining substantial acquisition.}
\textbf{a, Sources and continuation design.} The compact, bilobed and annular labels refer to predetermined structural origins, not three maintained morphologies or new species. Initial mass and stored energy are 24 M and 14.4 EU. The annular source uses different fixed physical-rule coefficients, so across-source differences do not isolate initial geometry. The displayed acquired states are the three updating stripe histories at completed step 512. In the first stripe phase the bilobed updating branch has $\Delta E=-3.671717$ EU relative to its paired actor-increment-off branch; it is retained in the continuation comparison. Initial images use $x=6\ldots27,y=6\ldots27$; acquired-state images use $x=8\ldots29,y=26\ldots31,0\ldots15$. Each row uses one fixed window across all origins, with no rotational or translational registration. White-to-teal density uses a common square-root display scale, 0--1.2 M/cell, with no saturated selected cells. Displayed windows retain at least 99.66\% of the corresponding full-domain mass in panel a; statistics use the full domain.

The complete updating stripe state at 512 is restored before assigning either continued conditions or the joint perturbation: remaining stock and its support shift by eight periodic $x$ cells without replenishment, and access is open for 16 actions per 64 through 768, then 16 per 128. Within each future, only further $\theta$ increments are enabled or disabled. The no-increment branch retains and transports/mixes its acquired spatial logits; it is not a zero-memory or mean-policy control. Critic, recent-food and original $m/z$ processes remain online. A new common random tape is shared by the four continuations of each acquired body. The run/brake actor has no interval-$T$ controller. Braking suppresses the autonomous polarity readout; transported polarity and morphology-driven motion persist. Three receiver histories produce twelve continuation branches, not twelve independent acquisition histories. The stripe subset was chosen after the first batch and does not represent all acquired states.

\textbf{b, A representative conditional process.} Compact, shifted/pulsed branches are shown at 512, 768 and 1024 in a common unregistered original-coordinate window, $x=6\ldots29,y=26\ldots31,0\ldots17$. Each image is one archived density field. The green native overlay marks the left edge of the resource support at $x=22$; the support extends to the right edge of the window ($x=30$). It is geometric support, not a claim of open access or active uptake. The selected windows retain at least 99.55\% of full-domain mass. The resource-support curve sums completed-state tissue over the fixed support, irrespective of instantaneous access. Both displayed branches start from $M=23.499519$ and $E=38.840443$; at 1024 the updating/no-increment branches have $M=23.461774/11.998454$ and $E=3.516099/0.021080$. Over actions 513--1024, updating receives 4.635672 EU and pays 39.958639 EU, while the no-increment branch receives 4.352202 EU and pays 43.074004 EU. Actual payment is whole-body motor plus basal payment, excluding energy removed with tissue. It is not a source-to-destination causal decomposition. Native curves are geometrically simplified displays of all 512 archived action rows plus the initial state; endpoint and cumulative totals use unsimplified values. Lower intake curves and upper payment curves use the same EU axis. The large mass contrast therefore does not imply abundant renewed supply: the updating branch consumes most of its initial reserve while contact and intake fall.

\textbf{c, All continuation contrasts.} Horizontal segments run from zero to the signed paired effect (updating minus no further actor increments); end ticks mark the values, with rounded signed labels for effects smaller than 0.3 in the respective units. Light-grey reference lines mark zero. At native drawing precision, subpixel effects can coincide with zero; signed labels preserve their direction. Grey denotes unchanged conditions and teal denotes the joint shift/access perturbation. All three energy columns use $[-4.5,4.5]$ EU; the mass column has its own M scale. No uncertainty interval or replicated-history estimate is implied. The small changed-bilobed effects ($\Delta M=0.686452$, $\Delta E=0.057131$ EU) and mixed unchanged-future outcomes are retained. These are conditional effects of continued updating in acquired states, not acquired-versus-naive transfer efficacy, proof of stable adaptation, or a mediation analysis. All data come from the completed 5031 and 96d1 batches (6144 steps each; 12288 total). Figure preparation adds zero model, body or energy updates.
